

%% 
%% This is a LaTeX document generated by automatic conversion of a
%% Mathematica notebook using Mathematica 4.0.
%% 
%% This document uses special macros that are defined in a LaTeX 
%% package file with a name of the form notebook*.sty.  Appropriate 
%% commands for loading the package have been included in this document.
%% The LaTeX package files can be found in a directory with path:
%% 
%% /usr/local/mathematica/SystemFiles/IncludeFiles/TeX/
%% 
%% The LaTeX package used in this document may require that your 
%% LaTeX implementation support fonts other than the standard Computer 
%% Modern fonts that are used by LaTeX.  See the Additional Information 
%% section under the TeXSave entry in the Mathematica Reference Guide 
%% for more information.
%% 
%% Implementation-specific instructions for installing TeX-related files
%% can be found at this URL.
%% 
%% http://www.wolfram.com/support/FrontEnds/Export/TeX
%% 
%% 

\documentclass{article}
\usepackage{notebook2e, latexsym}


\begin{document}

\Title{DISCRETE-TIME FOURIER SERIES}
\Subtitle{LECTURE 11}
\Subsubtitle{ELE314: LINEAR SIGNALS AND SYSTEMS}
\Subsubtitle{Mariusz Jankowski \copyright{}

University of Southern Maine

mjkcc@usm.maine.edu}
\dispSFinmath{<<\Mvariable{SignalPlot`}}
\begin{CellGroup}
\Section*{1 Introduction}


The techniques of discrete-time Fourier analysis are of great value in analyzing and gaining insight into properties of discrete-time signals and
systems.  In the 1940s and 1950s major steps forward were taken in the development of discrete-time techniques in general and in the use of the tools
of discrete-time Fourier analysis in particular.  In the mid-1960s an algorithm, now known as the Fast Fourier Transform or FFT, was developed (J.W.
Cooley and J.W. Tukey).  In this document the related decompositions: discrete-time Fourier series (DTFS), the Discrete Fourier Transform (DFT),
and the Fast Fourier Transform (FFT) will be presented.


\end{CellGroup}
\begin{CellGroup}
\Section*{2 Discrete-Time Sinusoidal Signals}


In analogy to the continuous time case, the basis functions of the DTFS are discrete-time complex exponentials. Recall from an earlier discussion
(see Lecture 2: \buttonbox[Hyperlink]{Basic One-Dimensional Signals}), that given a period {\itshape M}, the family of discrete-time complex exponentials
consists of all sinusoidal signals with frequencies that are integer multiples of the some fundamental frequency \inlineTFinmath{\frac{2\pi }{M}}.
 It has been shown that for any integer {\itshape M}, there is a finite number {\itshape M} of the basis functions.  As a consequence, the DTFS and
similarly the closely related Discrete Fourier Transform (DFT), are finite length expansions. Consider a plot of all sinusoids of period \inlineTFinmath{M=12}.



\dispSFinmath{\MathBegin{MathArray}{l}
\Mfunction{Partition}\big[\Mfunction{Table}\big[\Muserfunction{LinePlot}\big[\Mfunction{Table}\big[\cos \big[\frac{\pi \multsp n\multsp k}{4}\big],\multsp
\{n,-7,8\}\big],  \\
\noalign{\vspace{0.90625ex}}\hspace{1em}\hspace{1em}\hspace{1em}\hspace{1em}\hspace{1em}\hspace{1em}\Mvariable{PlotLabel}\rightarrow \Mvariable{frecuencia\multsp
}<>\Mfunction{ToString}[k],  \\
\noalign{\vspace{0.666667ex}}
\hspace{3.em} \Mvariable{PointStyle}->\Mfunction{PointSize}[0.04],\Mvariable{DisplayFunction}\rightarrow \Mvariable{Identity}\big],  \\
\noalign{\vspace{0.833333ex}}\hspace{1em}\hspace{1em}\{k,-7,8\}\big],2\big];\\
\MathEnd{MathArray}}
\begin{CellGroup}
\dispSFinmath{\Mfunction{Show}[\Mfunction{GraphicsArray}[\%]];}
\mathGraphic{lecture11.tex_gr1.eps}

\end{CellGroup}

\end{CellGroup}
\begin{CellGroup}
\Section*{3 Signal Decomposition and Reconstruction}


The discrete-time complex Fourier series is a representation of a periodic discrete-time signal (i.e. \inlineTFinmath{x[n]\multsp ==\multsp x[n+N]}
) in terms of linear combinations of harmonically related complex exponentials. The discrete-time Fourier series is defined 



\inlineTFinmath{x[n]\multsp =\multsp \multsp \sum _{k=0}^{N-1}X[k]\multsp \multsp {{\ExponentialE }^{\multsp j\multsp \frac{2\pi }{N}k\multsp n\multsp
}}},

for \inlineTFinmath{n=0,1,...,\multsp N-1}. \inlineTFinmath{X[k]} is the {\itshape k}-th Fourier coefficient of the periodic sequence \inlineTFinmath{x[n]}.
Note that the series defined by (1) differs in a fundamental and important way from the series formulation of continuous-time periodic signals -
it is finite! This is a consequence of the periodic nature of the Fourier coefficients (see Problem 1, below) and has far-reaching computational
consequences. The summation in (1) can  be defined over any interval of length {\itshape N}, for example \inlineTFinmath{k=l,l+1,...,l+N-1}.



The DTFS analysis equation plays the same role for discrete-time signals as it does for continuous-time signals. It provides a formula for the DTFS
coefficients \inlineTFinmath{X[k]} 's of a periodic signal.



\inlineTFinmath{X[k]\multsp =\multsp \frac{1}{N}\multsp \multsp \sum _{n=0}^{N-1}x[n]\multsp {{\ExponentialE }^{-\multsp j\multsp \frac{2\multsp
\pi }{N}\multsp \multsp k\multsp \multsp n}}}, 

where {\itshape N} is the period of the signal \inlineTFinmath{x[n]} ({\itshape N} is an integer) and \inlineTFinmath{k} is usually evaluated on
the finite integer interval, \inlineTFinmath{k=0,1,...,M-1}.



\begin{CellGroup}
\Exercise{Problem 1}
\ExerciseText{Prove that the FS coefficients of a periodic, discrete-time signal with period \inlineTFinmath{M}, are themselves periodic with period
\inlineTFinmath{M}, i.e. \inlineTFinmath{X[k+N]\multsp =\multsp X[k]} .}



\end{CellGroup}
Recall that in the continuous-time case, the analysis formula was defined by an integral over one period, whereas here the calculation requires a
summation over a finite set of complex exponentials. Following our earlier discussion of the continuous-time FS, we proceed with a simple example
of evaluating (2) for a discrete-time periodic square wave. We define a signal \inlineTFinmath{x[n]} , 



\inlineTFinmath{x[n]\multsp =\multsp \big\{\MathBegin{MathArray}[c]{cc}
	1,&\Mvariable{for}\multsp \multsp -4+l\multsp N\leq n\leq 4+l\multsp N \\
	0,&\Mvariable{else} 
  \MathEnd{MathArray}}

with \inlineTFinmath{N=32}, and \inlineTFinmath{l\multsp \epsilon \multsp \DoubleStruckCapitalZ }. 



\mathGraphic{lecture11.tex_gr2.eps}
Substituting (3) into (2) and simplifying, we get the following expression for the Fourier coefficients of the square wave.

\dispSFinmath{X[\Mvariable{k\_}]:=\frac{1}{32}\bigg(\sum _{n=0}^{4}\exp \big[-\frac{\imag \multsp \multsp \pi \multsp }{16}\multsp n\multsp k\big]\multsp
+\multsp \sum _{n=28}^{31}\exp \big[-\frac{\imag \multsp \multsp \pi \multsp }{16}\multsp n\multsp k\big]\multsp \bigg);}
\begin{CellGroup}
\dispSFinmath{\Muserfunction{LinePlot}[\Mfunction{Table}[\{k,X[k]\},\{k,0,31\}],\Mvariable{PlotRange}->\Mvariable{All}];}
\mathGraphic{lecture11.tex_gr3.eps}

\end{CellGroup}
Observe carefully the similarities between the continuous and discrete-time results. The discrete-time result has the general form of a sinc function,
however, as established in Problem 1, the set of FS coefficients is periodic with period \inlineTFinmath{M}. The coefficients of a real sequence
have conjugate symmetry (even symmetry for real valued coefficients). The coefficients of a discrete-time signal with even symmetry (as in our example)
are real valued (for a complete listing of properties of FS coefficients see [Oppenheim]).



\begin{CellGroup}
\Exercise{Problem 2}
\ExerciseText{Obtain the DTFS coefficients of the signal }
\begin{CellGroup}
\dispSFinmath{\cos \big[\frac{2\multsp \pi }{8}\multsp 3\multsp \multsp n\big]}
\dispSFoutmath{\cos \big[\frac{3\multsp n\multsp \pi }{4}\big]}

\end{CellGroup}
\ExerciseText{by Euler's theorem and by direct evaluation. (Hint: first determine the period {\itshape M} of the given signal). }



\end{CellGroup}
\begin{CellGroup}
\Exercise{Problem 3}
\ExerciseText{Obtain the period of the following signal: \inlineTFinmath{\cos\Big(\frac{2\pi }{7}3\multsp n\multsp \Big)}. Calculate the DTFS coefficients
of the signal using the actual period ({\itshape i.e.}, \inlineTFinmath{N=7}) and the value \inlineTFinmath{N=8}. Compare the two results. Try to
explain the results (this is NOT easy!).}



\end{CellGroup}
Again, in analogy with the continuous-time case, it is of interest to investigate the properties of signal reconstruction based on an incomplete
set of Fourier coefficients. In this particular example consider the result of computing \inlineTFinmath{\overvar{x}{\RawWedge }[n]} from the 5 lowest
order harmonics (i.e. \inlineTFinmath{k=0,\pm 1,...,\pm 5}).  Compare with the original signal \inlineTFinmath{x[n]}.  What do you notice?



\dispSFinmath{a[\Mvariable{n\_}]:=\multsp \sum _{k=-5}^{5}\multsp X[k]*\exp \big[\imag \frac{\pi }{16}\multsp \multsp n\multsp k\big]}
\begin{CellGroup}
\dispSFinmath{\Muserfunction{LinePlot}[\Mfunction{Table}[\{n,a[n]\},\{n,-16,40\}],\Mvariable{PointStyle}->\Mfunction{PointSize}[0.02]];}
\mathGraphic{lecture11.tex_gr4.eps}

\end{CellGroup}
Clearly, the reconstructed signal is no longer the ideal periodic pulse defined in (3).


\end{CellGroup}
\begin{CellGroup}
\Section*{4 DTFS as a System of Linear Equations}


Consider a discrete-time signal \inlineTFinmath{x[n]} of period M. From the synthesis formula (1) in the previous section, it is clear that the signal
can be represented in terms of the Fourier series. Instead of using the analysis formula (2) to obtain the Fourier coefficients, consider the synthesis
formula as a set of N linear equations:



\inlineTFinmath{\MathBegin{MathArray}[c]{c}
	x[0]=\sum _{k=0}^{N-1}X[k] \\
	x[1]=\sum _{k=0}^{N-1}X[k]\multsp \multsp {{\ExponentialE }^{\multsp j\multsp \frac{\multsp 2\multsp \pi \multsp }{N}k}} \\
	... \\
	x[N-1]=\sum _{k=0}^{N-1}X[k]\multsp {{\ExponentialE }^{\multsp j\multsp \frac{\multsp 2\multsp \pi \multsp }{N}(N-1)k}} 
  \MathEnd{MathArray}}

There are N unknowns, \inlineTFinmath{X[0],\multsp X[1],\multsp ...\multsp ,\multsp X[N-1]}, in this system of N independent linear equations, and
so they can be easily obtained by standard methods.



\begin{CellGroup}
\Exercise{Problem 4}
\ExerciseText{(a) Explicitly write all the equations for the case \inlineTFinmath{N=4}.

(b) Given one period of the signal \inlineTFinmath{x[n]}, solve for the Fourier series coefficients \inlineTFinmath{X[0],\multsp X[1],X[2],X[3]\multsp
} using standard methods of solving systems of linear equations (Solve or LinearSolve). Use a sequence \inlineTFinmath{x[n]} with the following values:}


\begin{CellGroup}
\dispSFinmath{x=\{1,0,2,-1\}}
\dispSFoutmath{\{1,0,2,-1\}}

\end{CellGroup}

\end{CellGroup}

\end{CellGroup}
\begin{CellGroup}
\Section*{PROBLEM SOLUTIONS}
\begin{CellGroup}
\Exercise{Problem 1}
\ExerciseText{Prove that the FS coefficients of a periodic, discrete-time signal with period \inlineTFinmath{M}, are themselves periodic with period
\inlineTFinmath{M}, i.e. \inlineTFinmath{X[k+M]\multsp =\multsp X[k]} .}



\end{CellGroup}
\begin{CellGroup}
\Exercise{Problem 2}
\ExerciseText{Obtain the DTFS coefficients of the signal }
\begin{CellGroup}
\dispSFinmath{\cos \big[\frac{2\multsp \pi }{8}\multsp 3\multsp \multsp n\big]}
\dispSFoutmath{\cos \big[\frac{3\multsp n\multsp \pi }{4}\big]}

\end{CellGroup}
\ExerciseText{by Euler's theorem and by direct evaluation. (Hint: first determine the period M of the given signal). }

\end{CellGroup}
Using the definition we get

\dispSFinmath{\Mfunction{Clear}[X]}
\dispSFinmath{X[\Mvariable{k\_}]:=\frac{1}{8}\sum _{n=0}^{7}\cos \big[\frac{2\multsp \pi }{8}\multsp 3\multsp \multsp n\big]\multsp \exp \big[-\imag
\multsp \frac{\pi }{4}k\multsp n\big]}
\begin{CellGroup}
\dispSFinmath{\Mfunction{Table}[\{k,X[k]\},\{k,0,7\}]}
\dispSFoutmath{\MathBegin{MathArray}{l}
\big\{\{0,0\},\big\{1,\frac{1}{8}\multsp \Bigg(2-\frac{{{\ExponentialE }^{-\frac{\ImaginaryI \multsp \pi }{4}}}}{{\sqrt{2}}}-\frac{{{\ExponentialE
}^{\frac{\ImaginaryI \multsp \pi }{4}}}}{{\sqrt{2}}}+\frac{{{\ExponentialE }^{-\frac{3\multsp \ImaginaryI \multsp \pi }{4}}}}{{\sqrt{2}}}+\frac{{{\ExponentialE
}^{\frac{3\multsp \ImaginaryI \multsp \pi }{4}}}}{{\sqrt{2}}}\Bigg)\big\},  \\
\noalign{\vspace{2.ex}}
\hspace{1.em} \{2,0\},\big\{3,\frac{1}{8}\multsp \Bigg(2+\frac{{{\ExponentialE }^{-\frac{\ImaginaryI \multsp \pi }{4}}}}{{\sqrt{2}}}+\frac{{{\ExponentialE
}^{\frac{\ImaginaryI \multsp \pi }{4}}}}{{\sqrt{2}}}-\frac{{{\ExponentialE }^{-\frac{3\multsp \ImaginaryI \multsp \pi }{4}}}}{{\sqrt{2}}}-\frac{{{\ExponentialE
}^{\frac{3\multsp \ImaginaryI \multsp \pi }{4}}}}{{\sqrt{2}}}\Bigg)\big\},  \\
\noalign{\vspace{2.ex}}
\hspace{1.em} \{4,0\},\big\{5,\frac{1}{8}\multsp \Bigg(2+\frac{{{\ExponentialE }^{-\frac{\ImaginaryI \multsp \pi }{4}}}}{{\sqrt{2}}}+\frac{{{\ExponentialE
}^{\frac{\ImaginaryI \multsp \pi }{4}}}}{{\sqrt{2}}}-\frac{{{\ExponentialE }^{-\frac{3\multsp \ImaginaryI \multsp \pi }{4}}}}{{\sqrt{2}}}-\frac{{{\ExponentialE
}^{\frac{3\multsp \ImaginaryI \multsp \pi }{4}}}}{{\sqrt{2}}}\Bigg)\big\},  \\
\noalign{\vspace{2.ex}}
\hspace{1.em} \{6,0\},\big\{7,\frac{1}{8}\multsp \Bigg(2-\frac{{{\ExponentialE }^{-\frac{\ImaginaryI \multsp \pi }{4}}}}{{\sqrt{2}}}-\frac{{{\ExponentialE
}^{\frac{\ImaginaryI \multsp \pi }{4}}}}{{\sqrt{2}}}+\frac{{{\ExponentialE }^{-\frac{3\multsp \ImaginaryI \multsp \pi }{4}}}}{{\sqrt{2}}}+\frac{{{\ExponentialE
}^{\frac{3\multsp \ImaginaryI \multsp \pi }{4}}}}{{\sqrt{2}}}\Bigg)\big\}\big\}\\
\MathEnd{MathArray}}

\end{CellGroup}
\begin{CellGroup}
\dispSFinmath{\%//\Mfunction{N}}
\dispSFoutmath{\MathBegin{MathArray}{l}
\{\{0.,0.\},\{1.,-1.38778\times {{10}^{-17}}+0.\multsp \ImaginaryI \},\{2.,0.\},\{3.,0.5\InvisibleSpace +0.\multsp \ImaginaryI \},  \\
\noalign{\vspace{0.604167ex}}
\hspace{1.em} \{4.,0.\},\{5.,0.5\InvisibleSpace +0.\multsp \ImaginaryI \},\{6.,0.\},\{7.,-1.38778\times {{10}^{-17}}+0.\multsp \ImaginaryI \}\}\\
\MathEnd{MathArray}}

\end{CellGroup}
\begin{CellGroup}
\dispSFinmath{\Muserfunction{LinePlot}[\%\%,\Mvariable{PlotRange}->\Mvariable{All}]}
\mathGraphic{lecture11.tex_gr5.eps}
\dispSFoutmath{\SkeletonIndicator \Mvariable{Graphics}\SkeletonIndicator }

\end{CellGroup}
\begin{CellGroup}
\Exercise{Problem 3}
\ExerciseText{Obtain the period of the following signal: \inlineTFinmath{\cos\Big(\frac{2\pi }{7}3\multsp n\multsp \Big)}. Calculate the DTFS coefficients
of the signal using a period \inlineTFinmath{M=8}. Compare with the result of Problem 2. Explain.}



\end{CellGroup}
\begin{CellGroup}
\Exercise{Problem 4}
\ExerciseText{(a) Explicitly write all the equations for the case M\(=\)4.

(b) Given one period of the signal \inlineTFinmath{x[n]}, solve for the Fourier series coefficients \inlineTFinmath{X[0],\multsp X[1],...,X[4]\multsp
} using standard methods of solving systems of linear equations. Use a sequence \inlineTFinmath{x[n]} with the following values:}


\begin{CellGroup}
\dispSFinmath{x=\{1,0,2,-1\}}
\dispSFoutmath{\{1,0,2,-1\}}

\end{CellGroup}

\end{CellGroup}

\end{CellGroup}
\begin{CellGroup}
\Section*{HOMEWORK PROBLEMS}
\begin{CellGroup}
\Subsubsection*{Homework Problem 11.1}
Determine the Fourier series coefficients of the following signals:

(a) Figure P3.28b (Oppenheim/Willsky, p. 258),

(b) Figure P3.28c (Oppenheim/Willsky, p. 258),

(c) \inlineTFinmath{x[n]} is periodic with period \inlineTFinmath{N=5}, and 

	\inlineTFinmath{x[n]\multsp =\multsp n}, for \inlineTFinmath{-2\leq n\leq 2}




\end{CellGroup}
\begin{CellGroup}
\Subsubsection*{Homework Problem 11.2}
Verify Parseval's relation using the discrete-time periodic square wave as an example. 

(a) Obtain the signal power \inlineTFinmath{{P_T}} using the time-domain calculation.

(b) Obtain the signal power \inlineTFinmath{{P_F}} by evaluating over the FS coefficients.

Show that \inlineTFinmath{{P_T}\multsp =\multsp {P_F}}.




\end{CellGroup}

\end{CellGroup}

\end{document}
